% Page 018 — page imprimée 14
% Suite de la section consacrée à Pierre Bertrand, galérien de Campis.


{\fontsize{10.55}{12.55}\selectfont

\vspace*{1.55cm}

\citationblock{%
« Comment une partie des forçats protestants ont-ils pu résister à la machine des
galères ? En construisant, peu à peu, une organisation clandestine de solidarité et de
résistance, qui exploita remarquablement toutes les failles de l’institution pénale. Un
texte de 1699, intitulé « Règlements faits sur les galères de France par les
confesseurs qui souffrent pour la vérité de l’Evangile », détaille le fonctionnement du
réseau des galériens protestants. Il était dirigé par une commission de 7 membres ;
Chacun d’eux s’occupait d’un ou plusieurs secteurs d’activité (propagande,
relations épistolaires, biographies de galériens pour la foi, distribution des secours,
contacts avec les protestants de France...) et devait périodiquement en rendre
compte à ses camarades.

L’entraide constituait le premier but. Au moyen des collectes organisées dans les
églises françaises de Suisse, d’Allemagne et de Hollande, où grâce à des actes de
générosité individuelle, de l’argent, et aussi des livres étaient régulièrement envoyés
aux forçats réformés. Il ne s’agissait pas d’un pactole mais plutôt de dons modestes
et de petites sommes... ». \textsuperscript{*}}

Ainsi les plus déterminés des galériens protestants se donnèrent clandestinement les
moyens de résister. En utilisant toutes les ressources, licites ou illicites, de la société
des galériens, ils purent s’entraider, prier, lire, écrire, s’instruire même...

\vspace{0.25em}

\citationblock{%
« Ces condamnés politiques » du Roi Soleil, n’ont bénéficié d’aucun régime de
faveur. Jusqu’à la conclusion de la paix d’Utrecht (1713) le roi ne gracie que ceux
qui se convertissent en galère. Les opiniâtres qui survécurent et refusèrent d’abjurer,
reçurent un ordre de liberté qui comportait une clause particulière. Leur libération
fut effective, en 1713 ou 1714, grâce aux pressions de la reine d’Angleterre (Anne
Stuart). Mais Louis XIV exigea, leur sortie du Royaume. La plupart partirent à
Genève.}

\vspace{0.55em}

\begin{center}
\fbox{%
\begin{minipage}{0.78\textwidth}
\centering
44\% (700 environ) des galériens protestants sont morts à la peine.
\end{minipage}}
\end{center}

\vspace{0.65em}

\begin{center}
\textbf{Pierre Bertrand de Campis, condamné aux galères (à vie) en 1698}\\
\textbf{est mort le 21 août 1705 à l’âge de 29 ans}\\
\textbf{à l’hôpital des forçats de Marseille}
\end{center}

\vfill

\begin{center}
\itshape
\textsuperscript{*} André Zysberg : \textit{Les Galériens : Vies et destins de 60 000
forçats sur les galères de France 1680 - 1748}
\end{center}

}

